\documentclass[12pt,a4paper]{article}
\usepackage[utf8]{inputenc}
\usepackage[T1]{fontenc}
\usepackage{lmodern}
\usepackage[margin=1in]{geometry}
\usepackage{enumitem}
\usepackage{fancyhdr}
\usepackage{xcolor}
\usepackage{tcolorbox}
\usepackage{amsmath}
\usepackage{amssymb}

% Header/Footer
\pagestyle{fancy}
\fancyhf{}
\rhead{Graph Embeddings \& Retrieval Assessment}
\lhead{100 Questions}
\cfoot{\thepage}

\definecolor{easy}{rgb}{0.2,0.6,0.2}
\definecolor{medium}{rgb}{0.8,0.6,0.0}
\definecolor{hard}{rgb}{0.8,0.2,0.2}

\title{\Huge\textbf{Graph Embeddings \& Retrieval}\\[0.3cm]
\LARGE Mastery Quiz\\[0.5cm]
\Large 100 Multiple Choice Questions\\[0.3cm]
\normalsize Progressive Difficulty: Easy → Expert}
\author{Self-Assessment for Machine Learning Engineers}
\date{January 2026}

\begin{document}

\maketitle
\thispagestyle{empty}

\vspace{1cm}

\begin{center}
\begin{tcolorbox}[width=0.8\textwidth,colback=gray!5,colframe=gray!50]
\textbf{Instructions:}
\begin{itemize}
    \item Questions 1-30: \textcolor{easy}{\textbf{Easy}} - Fundamentals
    \item Questions 31-60: \textcolor{medium}{\textbf{Medium}} - Applied concepts
    \item Questions 61-85: \textcolor{hard}{\textbf{Hard}} - Advanced topics
    \item Questions 86-100: \textcolor{hard}{\textbf{Expert}} - Research-level understanding
\end{itemize}
Each question has exactly ONE correct answer. Recommended time: 120 minutes.
\end{tcolorbox}
\end{center}

\newpage

%==============================================================================
% EASY QUESTIONS (1-30)
%==============================================================================
\section*{\textcolor{easy}{Section 1: Easy (Questions 1-30)}}

\begin{enumerate}

% Q1 - Answer: B
\item What is a vector in the context of machine learning?
\begin{enumerate}[label=\Alph*)]
    \item A type of database
    \item An ordered list of numbers representing features
    \item A graph structure
    \item A neural network layer
\end{enumerate}

% Q2 - Answer: D
\item What does the dot product of two vectors measure?
\begin{enumerate}[label=\Alph*)]
    \item The distance between them
    \item Their multiplication
    \item Their sum
    \item How aligned they are in direction
\end{enumerate}

% Q3 - Answer: A
\item In a graph $G = (V, E)$, what does $V$ represent?
\begin{enumerate}[label=\Alph*)]
    \item The set of nodes (vertices)
    \item The set of edges
    \item The vector space
    \item The adjacency matrix
\end{enumerate}

% Q4 - Answer: C
\item What is an edge in a graph?
\begin{enumerate}[label=\Alph*)]
    \item A node's feature vector
    \item A node's embedding
    \item A connection between two nodes
    \item The boundary of the graph
\end{enumerate}

% Q5 - Answer: B
\item What is cosine similarity used for?
\begin{enumerate}[label=\Alph*)]
    \item Counting graph edges
    \item Measuring similarity between vectors
    \item Training neural networks
    \item Computing graph density
\end{enumerate}

% Q6 - Answer: D
\item What is the range of cosine similarity values?
\begin{enumerate}[label=\Alph*)]
    \item 0 to infinity
    \item 0 to 1
    \item -infinity to infinity
    \item -1 to 1
\end{enumerate}

% Q7 - Answer: A
\item What is an embedding in machine learning?
\begin{enumerate}[label=\Alph*)]
    \item A dense vector representation of an object
    \item A type of graph
    \item A loss function
    \item An optimization algorithm
\end{enumerate}

% Q8 - Answer: C
\item What does the degree of a node represent?
\begin{enumerate}[label=\Alph*)]
    \item Its position in the graph
    \item Its embedding dimension
    \item The number of edges connected to it
    \item Its importance score
\end{enumerate}

% Q9 - Answer: B
\item What is an adjacency matrix?
\begin{enumerate}[label=\Alph*)]
    \item A matrix of node features
    \item A matrix showing which nodes are connected
    \item A matrix of embeddings
    \item A matrix of similarities
\end{enumerate}

% Q10 - Answer: D
\item In retrieval systems, what is a query?
\begin{enumerate}[label=\Alph*)]
    \item The database
    \item The index structure
    \item The retrieved documents
    \item The input used to search for relevant items
\end{enumerate}

% Q11 - Answer: A
\item What is TF-IDF primarily used for?
\begin{enumerate}[label=\Alph*)]
    \item Weighing word importance in documents
    \item Training graph neural networks
    \item Building adjacency matrices
    \item Computing node degrees
\end{enumerate}

% Q12 - Answer: C
\item What is a directed graph?
\begin{enumerate}[label=\Alph*)]
    \item A graph with weighted edges
    \item A graph with node features
    \item A graph where edges have direction (A→B doesn't imply B→A)
    \item A graph with multiple components
\end{enumerate}

% Q13 - Answer: B
\item What is the purpose of NumPy's \texttt{np.dot()} function?
\begin{enumerate}[label=\Alph*)]
    \item Create a new array
    \item Compute the dot product of arrays
    \item Plot data points
    \item Load files
\end{enumerate}

% Q14 - Answer: D
\item What does \texttt{np.linalg.norm()} compute?
\begin{enumerate}[label=\Alph*)]
    \item Matrix inverse
    \item Eigenvalues
    \item Matrix determinant
    \item Vector or matrix norm (magnitude)
\end{enumerate}

% Q15 - Answer: A
\item In graph notation, what does $\mathcal{N}(v)$ typically represent?
\begin{enumerate}[label=\Alph*)]
    \item The neighbors of node $v$
    \item The normal distribution
    \item The number of nodes
    \item The neural network
\end{enumerate}

% Q16 - Answer: C
\item What makes graph data different from tabular data?
\begin{enumerate}[label=\Alph*)]
    \item Graphs have more rows
    \item Graphs are always larger
    \item Graphs capture relationships between entities
    \item Graphs cannot have numerical features
\end{enumerate}

% Q17 - Answer: B
\item What is a random walk on a graph?
\begin{enumerate}[label=\Alph*)]
    \item Randomly generating a new graph
    \item A sequence of nodes where each step moves to a random neighbor
    \item Randomly removing edges
    \item A random initialization of embeddings
\end{enumerate}

% Q18 - Answer: D
\item What is FAISS used for?
\begin{enumerate}[label=\Alph*)]
    \item Training neural networks
    \item Graph visualization
    \item Text preprocessing
    \item Fast similarity search in high-dimensional vectors
\end{enumerate}

% Q19 - Answer: A
\item What is Precision@K in retrieval?
\begin{enumerate}[label=\Alph*)]
    \item Fraction of top-K results that are relevant
    \item Total number of relevant documents
    \item The K most precise documents
    \item Accuracy of the ranking model
\end{enumerate}

% Q20 - Answer: C
\item What is Recall@K in retrieval?
\begin{enumerate}[label=\Alph*)]
    \item How many results to return
    \item The precision of the model
    \item Fraction of all relevant documents found in top-K
    \item The number of queries processed
\end{enumerate}

% Q21 - Answer: B
\item What is PyTorch Geometric (PyG)?
\begin{enumerate}[label=\Alph*)]
    \item A visualization library
    \item A library for deep learning on graphs
    \item A database system
    \item A data loading utility
\end{enumerate}

% Q22 - Answer: D
\item What does L2 normalization do to a vector?
\begin{enumerate}[label=\Alph*)]
    \item Doubles its length
    \item Removes negative values
    \item Sorts its elements
    \item Scales it to have unit length (magnitude 1)
\end{enumerate}

% Q23 - Answer: A
\item What is the purpose of an activation function in neural networks?
\begin{enumerate}[label=\Alph*)]
    \item Introduce non-linearity
    \item Store weights
    \item Compute gradients
    \item Load data
\end{enumerate}

% Q24 - Answer: C
\item What does ReLU stand for?
\begin{enumerate}[label=\Alph*)]
    \item Recursive Linear Unit
    \item Regular Linear Update
    \item Rectified Linear Unit
    \item Random Linear Utility
\end{enumerate}

% Q25 - Answer: B
\item In supervised learning, what are labels?
\begin{enumerate}[label=\Alph*)]
    \item Feature names
    \item Ground truth outputs for training data
    \item Model parameters
    \item Data preprocessing steps
\end{enumerate}

% Q26 - Answer: D
\item What is gradient descent?
\begin{enumerate}[label=\Alph*)]
    \item A type of graph
    \item A normalization technique
    \item A random walk strategy
    \item An optimization algorithm to minimize loss
\end{enumerate}

% Q27 - Answer: A
\item What is the typical range of embedding dimensions?
\begin{enumerate}[label=\Alph*)]
    \item 32 to 1024
    \item 1 to 5
    \item 10000 to 100000
    \item Always exactly 100
\end{enumerate}

% Q28 - Answer: C
\item In semantic search, what does "semantic" refer to?
\begin{enumerate}[label=\Alph*)]
    \item Exact keyword matching
    \item Speed of retrieval
    \item Meaning and context
    \item Database structure
\end{enumerate}

% Q29 - Answer: B
\item What is a citation network?
\begin{enumerate}[label=\Alph*)]
    \item A network of web pages
    \item A graph where nodes are papers and edges are citations
    \item A social media network
    \item A network of keywords
\end{enumerate}

% Q30 - Answer: D
\item What does \texttt{model.eval()} do in PyTorch?
\begin{enumerate}[label=\Alph*)]
    \item Evaluates the loss
    \item Saves the model
    \item Trains the model
    \item Sets the model to evaluation mode (disables dropout, etc.)
\end{enumerate}

%==============================================================================
% MEDIUM QUESTIONS (31-60)
%==============================================================================
\newpage
\section*{\textcolor{medium}{Section 2: Medium (Questions 31-60)}}

% Q31 - Answer: C
\item What is the key insight behind DeepWalk?
\begin{enumerate}[label=\Alph*)]
    \item Use supervised learning on graphs
    \item Compute exact graph distances
    \item Treat random walks as sentences and apply Word2Vec
    \item Directly minimize node distance
\end{enumerate}

% Q32 - Answer: A
\item In Node2Vec, what does the return parameter $p$ control?
\begin{enumerate}[label=\Alph*)]
    \item Likelihood of immediately returning to the previous node
    \item The embedding dimension
    \item The number of walks
    \item The learning rate
\end{enumerate}

% Q33 - Answer: D
\item What is message passing in GNNs?
\begin{enumerate}[label=\Alph*)]
    \item Sending gradients to neighbors
    \item Text communication between nodes
    \item Loading data in batches
    \item Aggregating information from neighboring nodes
\end{enumerate}

% Q34 - Answer: B
\item In the GCN formula, why do we add self-loops ($\tilde{A} = A + I$)?
\begin{enumerate}[label=\Alph*)]
    \item To reduce computation
    \item To include a node's own features in its update
    \item To make the matrix invertible
    \item To enable GPU computation
\end{enumerate}

% Q35 - Answer: C
\item What is the purpose of the degree normalization in GCN?
\begin{enumerate}[label=\Alph*)]
    \item Speed up training
    \item Reduce memory usage
    \item Prevent high-degree nodes from dominating
    \item Enable parallel computation
\end{enumerate}

% Q36 - Answer: A
\item How does GraphSAGE differ from GCN?
\begin{enumerate}[label=\Alph*)]
    \item GraphSAGE samples neighbors instead of using all
    \item GraphSAGE only works on directed graphs
    \item GraphSAGE doesn't use neural networks
    \item GCN is more scalable than GraphSAGE
\end{enumerate}

% Q37 - Answer: D
\item What does GAT learn that GCN does not?
\begin{enumerate}[label=\Alph*)]
    \item Node features
    \item Graph structure
    \item Embedding dimensions
    \item Attention weights for different neighbors
\end{enumerate}

% Q38 - Answer: B
\item In contrastive learning, what is a positive pair?
\begin{enumerate}[label=\Alph*)]
    \item Two unrelated samples
    \item Two samples that should be similar
    \item The model's predictions
    \item The training labels
\end{enumerate}

% Q39 - Answer: C
\item What is the purpose of temperature in InfoNCE loss?
\begin{enumerate}[label=\Alph*)]
    \item Control learning rate
    \item Normalize embeddings
    \item Control the sharpness of the softmax distribution
    \item Set batch size
\end{enumerate}

% Q40 - Answer: A
\item What does triplet loss optimize?
\begin{enumerate}[label=\Alph*)]
    \item Pushes positives closer and negatives farther from anchor
    \item Classifies samples into categories
    \item Reconstructs input data
    \item Minimizes reconstruction error
\end{enumerate}

% Q41 - Answer: D
\item What is the advantage of IVF (Inverted File) indexing over brute force?
\begin{enumerate}[label=\Alph*)]
    \item Higher recall
    \item Exact results
    \item Smaller index size
    \item Faster search by clustering similar vectors
\end{enumerate}

% Q42 - Answer: B
\item What is HNSW (Hierarchical Navigable Small World)?
\begin{enumerate}[label=\Alph*)]
    \item A type of GNN
    \item A graph-based approximate nearest neighbor algorithm
    \item A data augmentation technique
    \item A loss function
\end{enumerate}

% Q43 - Answer: C
\item Why normalize embeddings before cosine similarity search?
\begin{enumerate}[label=\Alph*)]
    \item To reduce memory
    \item To speed up training
    \item So inner product equals cosine similarity
    \item To enable backpropagation
\end{enumerate}

% Q44 - Answer: A
\item What is Mean Average Precision (MAP)?
\begin{enumerate}[label=\Alph*)]
    \item Average of average precision across all queries
    \item Maximum precision at any rank
    \item Minimum precision threshold
    \item The mean of all precisions
\end{enumerate}

% Q45 - Answer: D
\item What is MRR (Mean Reciprocal Rank)?
\begin{enumerate}[label=\Alph*)]
    \item The mean of all ranks
    \item The rank of the most relevant document
    \item The minimum rank achieved
    \item Average of 1/rank\_of\_first\_relevant across queries
\end{enumerate}

% Q46 - Answer: B
\item In NDCG, what does the "G" stand for?
\begin{enumerate}[label=\Alph*)]
    \item Graph
    \item Gain
    \item Gradient
    \item Global
\end{enumerate}

% Q47 - Answer: C
\item Why does NDCG use a logarithmic discount?
\begin{enumerate}[label=\Alph*)]
    \item To normalize scores
    \item To handle negative values
    \item To penalize relevant items appearing lower in the ranking
    \item To speed up computation
\end{enumerate}

% Q48 - Answer: A
\item What is the purpose of dropout in neural networks?
\begin{enumerate}[label=\Alph*)]
    \item Prevent overfitting by randomly zeroing neurons during training
    \item Speed up forward pass
    \item Reduce model size
    \item Initialize weights
\end{enumerate}

% Q49 - Answer: D
\item What is the difference between inductive and transductive learning?
\begin{enumerate}[label=\Alph*)]
    \item Speed of training
    \item Type of loss function
    \item Number of layers
    \item Inductive generalizes to unseen nodes; transductive only works on seen nodes
\end{enumerate}

% Q50 - Answer: B
\item What does \texttt{torch.no\_grad()} do?
\begin{enumerate}[label=\Alph*)]
    \item Removes all gradients permanently
    \item Disables gradient computation for efficiency
    \item Freezes all model parameters
    \item Enables evaluation mode
\end{enumerate}

% Q51 - Answer: C
\item In PyTorch Geometric, what is \texttt{edge\_index}?
\begin{enumerate}[label=\Alph*)]
    \item A list of edge features
    \item A dictionary of neighbors
    \item A [2, num\_edges] tensor defining source and target nodes
    \item The adjacency matrix
\end{enumerate}

% Q52 - Answer: A
\item What is the purpose of a learning rate scheduler?
\begin{enumerate}[label=\Alph*)]
    \item Adjust learning rate during training based on progress
    \item Schedule when to save checkpoints
    \item Determine batch order
    \item Control GPU memory usage
\end{enumerate}

% Q53 - Answer: D
\item What is graph pooling used for?
\begin{enumerate}[label=\Alph*)]
    \item Adding new nodes
    \item Connecting disconnected components
    \item Removing noise
    \item Aggregating node embeddings into a graph-level embedding
\end{enumerate}

% Q54 - Answer: B
\item What is the advantage of mini-batch training for GNNs?
\begin{enumerate}[label=\Alph*)]
    \item Perfect accuracy
    \item Can handle graphs larger than GPU memory
    \item Eliminates the need for validation
    \item Guarantees convergence
\end{enumerate}

% Q55 - Answer: C
\item In Node2Vec, what does a low $q$ value encourage?
\begin{enumerate}[label=\Alph*)]
    \item Returning to previous node
    \item Staying at current node
    \item Exploring outward (BFS-like behavior)
    \item Shorter walks
\end{enumerate}

% Q56 - Answer: A
\item What is the receptive field in a GNN?
\begin{enumerate}[label=\Alph*)]
    \item The set of nodes whose features influence a node after L layers
    \item The size of the input features
    \item The output dimension
    \item The number of attention heads
\end{enumerate}

% Q57 - Answer: D
\item Why might deeper GNNs perform worse (over-smoothing)?
\begin{enumerate}[label=\Alph*)]
    \item Too many parameters
    \item Gradient vanishing
    \item Computational cost
    \item Node embeddings become too similar as they aggregate from larger neighborhoods
\end{enumerate}

% Q58 - Answer: B
\item What is the purpose of skip connections (residual connections)?
\begin{enumerate}[label=\Alph*)]
    \item Skip certain nodes during training
    \item Help gradients flow and preserve original features
    \item Skip validation steps
    \item Reduce computation
\end{enumerate}

% Q59 - Answer: C
\item What does a dual encoder architecture typically have?
\begin{enumerate}[label=\Alph*)]
    \item Two loss functions
    \item Two optimizers
    \item Separate encoders for queries and documents
    \item Two output layers
\end{enumerate}

% Q60 - Answer: A
\item What is negative sampling in contrastive learning?
\begin{enumerate}[label=\Alph*)]
    \item Selecting dissimilar samples to contrast with positives
    \item Removing outliers
    \item Downsampling the dataset
    \item Reducing batch size
\end{enumerate}

%==============================================================================
% HARD QUESTIONS (61-85)
%==============================================================================
\newpage
\section*{\textcolor{hard}{Section 3: Hard (Questions 61-85)}}

% Q61 - Answer: D
\item In GCN, why is the normalized Laplacian $D^{-1/2}AD^{-1/2}$ preferred over $D^{-1}A$?
\begin{enumerate}[label=\Alph*)]
    \item It's faster to compute
    \item It has fewer non-zero entries
    \item It requires less memory
    \item It's symmetric, leading to more stable eigenvalues and training
\end{enumerate}

% Q62 - Answer: B
\item What is the time complexity of a single GCN layer with $n$ nodes, $m$ edges, input dim $d$, output dim $d'$?
\begin{enumerate}[label=\Alph*)]
    \item $O(n^2)$
    \item $O(md + nd')$
    \item $O(n^3)$
    \item $O(d \cdot d')$
\end{enumerate}

% Q63 - Answer: C
\item In attention mechanisms, why do we divide by $\sqrt{d_k}$?
\begin{enumerate}[label=\Alph*)]
    \item To ensure values are between 0 and 1
    \item To normalize by number of heads
    \item To prevent dot products from becoming too large with high dimensions
    \item To match the learning rate
\end{enumerate}

% Q64 - Answer: A
\item What is the expressiveness limitation of 1-WL (Weisfeiler-Lehman) test?
\begin{enumerate}[label=\Alph*)]
    \item Cannot distinguish certain non-isomorphic graphs
    \item Cannot handle directed graphs
    \item Cannot process node features
    \item Cannot scale to large graphs
\end{enumerate}

% Q65 - Answer: D
\item Why is message passing GNNs' expressiveness bounded by 1-WL?
\begin{enumerate}[label=\Alph*)]
    \item They use the same number of layers
    \item They have the same computational complexity
    \item They use similar aggregation functions
    \item Both iteratively aggregate neighbor information in the same way
\end{enumerate}

% Q66 - Answer: B
\item What is spectral graph convolution based on?
\begin{enumerate}[label=\Alph*)]
    \item Random walks
    \item Graph Fourier transform using eigenvalues of Laplacian
    \item Attention mechanisms
    \item Neighborhood sampling
\end{enumerate}

% Q67 - Answer: C
\item What is the purpose of multi-head attention in GAT?
\begin{enumerate}[label=\Alph*)]
    \item Reduce computation
    \item Handle multiple graphs
    \item Learn different types of attention patterns in parallel
    \item Process multiple node types
\end{enumerate}

% Q68 - Answer: A
\item In approximate nearest neighbor search, what is recall@K?
\begin{enumerate}[label=\Alph*)]
    \item Fraction of true K nearest neighbors found
    \item Speed of retrieval
    \item Memory usage
    \item Index build time
\end{enumerate}

% Q69 - Answer: D
\item What is product quantization in vector search?
\begin{enumerate}[label=\Alph*)]
    \item Multiplying vectors before search
    \item A type of normalization
    \item Computing inner products
    \item Compressing vectors by splitting into subspaces and quantizing each
\end{enumerate}

% Q70 - Answer: B
\item In HNSW, what is the purpose of the hierarchical structure?
\begin{enumerate}[label=\Alph*)]
    \item Store more vectors
    \item Enable fast navigation from coarse to fine granularity
    \item Reduce memory usage
    \item Handle dynamic insertions
\end{enumerate}

% Q71 - Answer: C
\item What is the inductive bias of GNNs?
\begin{enumerate}[label=\Alph*)]
    \item Sequences matter
    \item Spatial locality matters
    \item Node features should be aggregated from neighbors
    \item Time ordering matters
\end{enumerate}

% Q72 - Answer: A
\item What is homophily in graphs?
\begin{enumerate}[label=\Alph*)]
    \item Connected nodes tend to have similar labels/features
    \item All nodes have the same degree
    \item The graph is fully connected
    \item Edges have uniform weights
\end{enumerate}

% Q73 - Answer: D
\item How do GNNs typically perform on heterophilic graphs?
\begin{enumerate}[label=\Alph*)]
    \item Always outperform other methods
    \item Same as on homophilic graphs
    \item Cannot process them
    \item Often worse, as aggregating dissimilar neighbors can hurt
\end{enumerate}

% Q74 - Answer: B
\item What is edge dropout in GNN training?
\begin{enumerate}[label=\Alph*)]
    \item Removing all edges
    \item Randomly removing edges during training for regularization
    \item Removing disconnected nodes
    \item Pruning the final graph
\end{enumerate}

% Q75 - Answer: C
\item In knowledge graph embeddings, what does TransE model?
\begin{enumerate}[label=\Alph*)]
    \item Node classification
    \item Graph generation
    \item Relations as translations: $h + r \approx t$
    \item Community detection
\end{enumerate}

% Q76 - Answer: A
\item What is link prediction?
\begin{enumerate}[label=\Alph*)]
    \item Predicting whether an edge should exist between two nodes
    \item Predicting node labels
    \item Predicting graph size
    \item Predicting the next random walk step
\end{enumerate}

% Q77 - Answer: D
\item In self-supervised learning on graphs, what is graph contrastive learning?
\begin{enumerate}[label=\Alph*)]
    \item Comparing graphs of different sizes
    \item Using human labels
    \item Minimizing reconstruction loss
    \item Creating augmented views and maximizing agreement between them
\end{enumerate}

% Q78 - Answer: B
\item What is the purpose of negative log-likelihood in training?
\begin{enumerate}[label=\Alph*)]
    \item Generate negative samples
    \item Maximize the probability of correct predictions
    \item Remove outliers
    \item Normalize embeddings
\end{enumerate}

% Q79 - Answer: C
\item Why might curriculum learning help in graph retrieval?
\begin{enumerate}[label=\Alph*)]
    \item Reduces memory usage
    \item Speeds up indexing
    \item Training on easy examples first helps the model learn better representations
    \item Enables multi-GPU training
\end{enumerate}

% Q80 - Answer: A
\item What is the "neighborhood explosion" problem in GNNs?
\begin{enumerate}[label=\Alph*)]
    \item The number of nodes to aggregate grows exponentially with depth
    \item Gradients become too large
    \item Memory usage increases linearly
    \item Training becomes unstable
\end{enumerate}

% Q81 - Answer: D
\item What technique does GraphSAGE use to handle neighborhood explosion?
\begin{enumerate}[label=\Alph*)]
    \item Attention mechanisms
    \item Skip connections
    \item Deeper networks
    \item Uniform sampling of a fixed number of neighbors
\end{enumerate}

% Q82 - Answer: B
\item What is the computational bottleneck in training retrievers with in-batch negatives?
\begin{enumerate}[label=\Alph*)]
    \item Forward pass
    \item Computing all-pairs similarity matrix
    \item Backpropagation
    \item Data loading
\end{enumerate}

% Q83 - Answer: C
\item In dense retrieval, what is the "embedding space collapse" problem?
\begin{enumerate}[label=\Alph*)]
    \item Embeddings become too large
    \item Memory runs out
    \item All embeddings converge to similar representations, losing discriminative power
    \item Training diverges
\end{enumerate}

% Q84 - Answer: A
\item What is the purpose of hard negative mining?
\begin{enumerate}[label=\Alph*)]
    \item Finding difficult negatives that are close to positives to improve learning
    \item Removing easy examples
    \item Balancing the dataset
    \item Reducing computation
\end{enumerate}

% Q85 - Answer: D
\item What is the asymptotic complexity of exact nearest neighbor search?
\begin{enumerate}[label=\Alph*)]
    \item $O(\log n)$
    \item $O(\sqrt{n})$
    \item $O(n \log n)$
    \item $O(n)$ per query
\end{enumerate}

%==============================================================================
% EXPERT QUESTIONS (86-100)
%==============================================================================
\newpage
\section*{\textcolor{hard}{Section 4: Expert (Questions 86-100)}}

% Q86 - Answer: C
\item Consider a 3-layer GCN on a graph where all nodes have identical features. What happens to node embeddings?
\begin{enumerate}[label=\Alph*)]
    \item Each node gets unique embeddings based on ID
    \item Embeddings depend only on node degree
    \item Nodes with the same 3-hop neighborhood structure get identical embeddings
    \item All nodes get identical embeddings regardless of structure
\end{enumerate}

% Q87 - Answer: B
\item In the attention mechanism $\alpha_{ij} = \text{softmax}_j(e_{ij})$, if we replace softmax with a learnable function that doesn't sum to 1, what property is lost?
\begin{enumerate}[label=\Alph*)]
    \item Gradient flow
    \item Interpretability as probability distribution and consistent scale
    \item Ability to have different weights
    \item Non-linearity
\end{enumerate}

% Q88 - Answer: A
\item What is the theoretical advantage of using higher-order WL tests (k-WL)?
\begin{enumerate}[label=\Alph*)]
    \item Can distinguish more non-isomorphic graphs
    \item Faster computation
    \item Lower memory usage
    \item Simpler implementation
\end{enumerate}

% Q89 - Answer: D
\item In InfoNCE loss with batch size $B$, how many negative samples does each query have?
\begin{enumerate}[label=\Alph*)]
    \item 1
    \item $B$
    \item $2B$
    \item $B-1$
\end{enumerate}

% Q90 - Answer: C
\item Consider two retrieval systems: one uses $\ell_2$ distance, the other uses inner product on L2-normalized vectors. Are they equivalent for ranking?
\begin{enumerate}[label=\Alph*)]
    \item No, they can produce different rankings
    \item Only for certain vector dimensions
    \item Yes, $\|\mathbf{a}-\mathbf{b}\|^2 = 2 - 2\mathbf{a}^T\mathbf{b}$ for unit vectors, so rankings are identical
    \item Only if all vectors have the same norm
\end{enumerate}

% Q91 - Answer: B
\item In the context of GNN expressiveness, what can higher-order message passing (k-GNNs) express that standard MPNNs cannot?
\begin{enumerate}[label=\Alph*)]
    \item Faster computation
    \item Subgraph patterns and higher-order structural features
    \item Larger graphs
    \item More node features
\end{enumerate}

% Q92 - Answer: A
\item What is the gradient of cosine similarity $\frac{\mathbf{a}^T\mathbf{b}}{\|\mathbf{a}\|\|\mathbf{b}\|}$ with respect to $\mathbf{a}$?
\begin{enumerate}[label=\Alph*)]
    \item $\frac{1}{\|\mathbf{a}\|}\left(\frac{\mathbf{b}}{\|\mathbf{b}\|} - \frac{\mathbf{a}^T\mathbf{b}}{\|\mathbf{a}\|\|\mathbf{b}\|}\frac{\mathbf{a}}{\|\mathbf{a}\|}\right)$
    \item $\frac{\mathbf{b}}{\|\mathbf{b}\|}$
    \item $\mathbf{b} - \mathbf{a}$
    \item $\frac{\mathbf{a} + \mathbf{b}}{2}$
\end{enumerate}

% Q93 - Answer: D
\item What is the relationship between margin in triplet loss and temperature in InfoNCE?
\begin{enumerate}[label=\Alph*)]
    \item They are identical
    \item Higher margin equals higher temperature
    \item They control different aspects
    \item Both control the "hardness" of learning: lower temperature / higher margin creates sharper decision boundaries
\end{enumerate}

% Q94 - Answer: C
\item In FAISS, what is the trade-off when increasing nprobe in IVF indexes?
\begin{enumerate}[label=\Alph*)]
    \item Higher nprobe = faster search, lower recall
    \item Higher nprobe = smaller index, higher recall
    \item Higher nprobe = higher recall, slower search
    \item nprobe doesn't affect recall
\end{enumerate}

% Q95 - Answer: B
\item What is the memory complexity of storing $n$ vectors of dimension $d$ as float32?
\begin{enumerate}[label=\Alph*)]
    \item $O(n)$
    \item $O(nd)$ bytes $= 4nd$ bytes
    \item $O(d)$
    \item $O(n^2d)$
\end{enumerate}

% Q96 - Answer: A
\item Consider training a retriever with cross-entropy over a similarity matrix. If batch size doubles, how does the gradient magnitude approximately change?
\begin{enumerate}[label=\Alph*)]
    \item Halves (since more negatives dilute the positive gradient)
    \item Doubles
    \item Stays the same
    \item Quadruples
\end{enumerate}

% Q97 - Answer: D
\item In graph attention, the attention weight $\alpha_{ij} = \frac{\exp(e_{ij})}{\sum_{k \in \mathcal{N}(i)} \exp(e_{ik})}$. If node $i$ has very high degree, what happens to individual $\alpha_{ij}$?
\begin{enumerate}[label=\Alph*)]
    \item They all become 1
    \item They become negative
    \item They become very large
    \item They become very small (each neighbor gets little attention)
\end{enumerate}

% Q98 - Answer: C
\item What is the issue with using random negatives vs. hard negatives early in training?
\begin{enumerate}[label=\Alph*)]
    \item Random negatives are computationally expensive
    \item Hard negatives are easier to find
    \item Hard negatives early on can destabilize training when the model can't distinguish them
    \item Random negatives lead to overfitting
\end{enumerate}

% Q99 - Answer: B
\item In a knowledge graph with entities $E$, relations $R$, and triples $T$, what is the parameter count of a TransE model with dimension $d$?
\begin{enumerate}[label=\Alph*)]
    \item $|T| \cdot d$
    \item $(|E| + |R|) \cdot d$
    \item $|E| \cdot |R| \cdot d$
    \item $|E|^2 \cdot d$
\end{enumerate}

% Q100 - Answer: A
\item Consider a retrieval system using a 2-layer GNN with $k$ sampled neighbors per layer. To compute an embedding for one query node, how many nodes need to be aggregated at most?
\begin{enumerate}[label=\Alph*)]
    \item $1 + k + k^2$ (exponential in depth)
    \item $2k$
    \item $k^2$
    \item $2^k$
\end{enumerate}

\end{enumerate}

%==============================================================================
% ANSWER KEY
%==============================================================================
\newpage
\section*{Answer Key}

\begin{center}
\textbf{Tear off or cover this page while taking the quiz.}
\end{center}

\vspace{1cm}

\begin{center}
\begin{tabular}{|c|c||c|c||c|c||c|c|}
\hline
\textbf{Q} & \textbf{Ans} & \textbf{Q} & \textbf{Ans} & \textbf{Q} & \textbf{Ans} & \textbf{Q} & \textbf{Ans} \\
\hline
1 & B & 26 & D & 51 & C & 76 & A \\
2 & D & 27 & A & 52 & A & 77 & D \\
3 & A & 28 & C & 53 & D & 78 & B \\
4 & C & 29 & B & 54 & B & 79 & C \\
5 & B & 30 & D & 55 & C & 80 & A \\
6 & D & 31 & C & 56 & A & 81 & D \\
7 & A & 32 & A & 57 & D & 82 & B \\
8 & C & 33 & D & 58 & B & 83 & C \\
9 & B & 34 & B & 59 & C & 84 & A \\
10 & D & 35 & C & 60 & A & 85 & D \\
11 & A & 36 & A & 61 & D & 86 & C \\
12 & C & 37 & D & 62 & B & 87 & B \\
13 & B & 38 & B & 63 & C & 88 & A \\
14 & D & 39 & C & 64 & A & 89 & D \\
15 & A & 40 & A & 65 & D & 90 & C \\
16 & C & 41 & D & 66 & B & 91 & B \\
17 & B & 42 & B & 67 & C & 92 & A \\
18 & D & 43 & C & 68 & A & 93 & D \\
19 & A & 44 & A & 69 & D & 94 & C \\
20 & C & 45 & D & 70 & B & 95 & B \\
21 & B & 46 & B & 71 & C & 96 & A \\
22 & D & 47 & C & 72 & A & 97 & D \\
23 & A & 48 & A & 73 & D & 98 & C \\
24 & C & 49 & D & 74 & B & 99 & B \\
25 & B & 50 & B & 75 & C & 100 & A \\
\hline
\end{tabular}
\end{center}

\vspace{1cm}

\subsection*{Scoring Guide}
\begin{itemize}
    \item \textbf{90-100}: Expert - Ready for research in graph ML
    \item \textbf{75-89}: Advanced - Strong foundation, review edge cases
    \item \textbf{60-74}: Intermediate - Solid progress, deepen mathematical understanding
    \item \textbf{45-59}: Beginner+ - Revisit core concepts
    \item \textbf{Below 45}: Beginner - Work through the tutorial again
\end{itemize}

\subsection*{Answer Distribution}
\begin{tabular}{|c|c|}
\hline
\textbf{Answer} & \textbf{Count} \\
\hline
A & 26 \\
B & 24 \\
C & 25 \\
D & 25 \\
\hline
\end{tabular}

\vspace{0.5cm}
Answers are distributed evenly and randomized---no patterns!

\subsection*{Topic Coverage}
\begin{itemize}
    \item Questions 1-10: Vectors, graphs, and embeddings basics
    \item Questions 11-20: Graph representations and retrieval basics
    \item Questions 21-30: PyTorch and neural network fundamentals
    \item Questions 31-45: Random walks, GNN architectures, loss functions
    \item Questions 46-60: Metrics, training, and architecture details
    \item Questions 61-75: Advanced GNN theory, search algorithms, expressiveness
    \item Questions 76-85: Link prediction, self-supervision, scaling
    \item Questions 86-100: Research-level theory, mathematical derivations, edge cases
\end{itemize}

\end{document}
